\documentclass[12pt]{article}
\usepackage{amsmath,amssymb,amsthm}

\newtheorem{theorem}{Theorem}
\newtheorem{lemma}[theorem]{Lemma}
\newtheorem{remark}[theorem]{Remark}

\begin{document}

\section*{Pullback Lemma for $C_P$ (Pinch(ii)) Reduction}

This document provides a rigorous 2-vertex cut pullback for the $C_P$ reduction.
We rely on the class $\mathcal{Q}$ of 3-connected cubic bipartite planar graphs and the definitions of a certified occurrence $H$ and its boundary $B$.

\begin{lemma}[Connectedness of $G \setminus H$]
\label{lem:isolated_neighborhood}
Let $G \in \mathcal{Q}$ contain a certified occurrence $H$. Then the induced subgraph $G[\,V(G)\setminus V(H)\,]$ is connected.
\end{lemma}
\begin{proof}
Since $G$ is 3-connected, it is 3-edge-connected ($\lambda(G) \ge 3$). Any component $C$ of $G[V(G)\setminus V(H)]$ would require $|E(C, V(H))| \ge 3$. If there were $k \ge 2$ such components, we would have $|E(V(G)\setminus V(H), V(H))| = \sum |E(C_i, V(H))| \ge 3k \ge 6$, contradicting the fact that a certified occurrence has exactly 4 boundary edges.
\end{proof}

\begin{lemma}[2-Vertex Cut Pullback for $C_P$]
\label{lem:pullback_cut_CP}
Assume $G \in \mathcal{Q}$ contains a certified $C_P$ occurrence $H$ with boundary $B = \{r, s, u_2, u_4\}$. Let $G' = (G \setminus H) \cup H'$ be the reduced graph, where $H'$ has adjacent gadget vertices $x, y$ such that $N_{G'}(x) \setminus \{y\} = \{r, s\} = B_x$ and $N_{G'}(y) \setminus \{x\} = \{u_2, u_4\} = B_y$.

If $G'$ has a 2-vertex cut $A = \{p, q\}$, then either:
\begin{enumerate}
    \item[(i)] $A \subseteq V(G) \setminus V(H)$ and $A$ is a 2-vertex cut of $G$; or
    \item[(ii)] exactly one of $\{p, q\}$ is a gadget vertex (w.l.o.g. $p=x$ or $p=y$) and $q \in V(G) \setminus V(H)$, and there exists a nonempty vertex set $K \subseteq V(G) \setminus V(H)$ such that every edge leaving $K$ in $G$ has its endpoint in $\{q\} \cup B_p$.
\end{enumerate}
\end{lemma}

\begin{proof}
Let $A = \{p, q\}$ be a 2-vertex cut in $G'$.

\medskip
\noindent\textbf{Case 1: $p, q \notin \{x, y\}$.} \\
Since the gadget edge $xy$ exists in $G'$, both $x$ and $y$ lie in the same component of $G' - A$; denote it $C_0$. Let $C_1$ be another component of $G' - A$.
Then $V(C_1) \subseteq V(G') \setminus \{x, y\} = V(G) \setminus V(H)$.
Any boundary terminal $b \in B \setminus A$ is adjacent in $G'$ to either $x$ or $y$ (both in $C_0$), so $b \in C_0$. Thus $C_1 \cap B = \emptyset$.
Restoring $H$ in $G$ does not add edges between $C_1$ and $V(H)$ because $C_1 \cap B = \emptyset$. Since internal edges of $V(G) \setminus V(H)$ are preserved, every edge in $G$ leaving $C_1$ must end in $A$. Thus $A$ is a 2-vertex cut in $G$, proving (i).

\medskip
\noindent\textbf{Case 2: $\{p, q\} = \{x, y\}$.} \\
As $G' - \{x, y\} = G[\,V(G) \setminus V(H)\,]$, Lemma~\ref{lem:isolated_neighborhood} implies this graph is connected. Thus $\{x, y\}$ cannot be a 2-vertex cut.

\medskip
\noindent\textbf{Case 3: Exactly one of $\{p, q\}$ lies in $\{x, y\}$.} \\
By symmetry, assume $p = x$ and $q \in V(G) \setminus V(H)$.
Let $C_0$ be the component of $G' - \{x, q\}$ containing $y$. Let $C_1$ be another component.
Define $K = V(C_1) \setminus B_x$. 
Note $V(C_1) \cap \{y\} = \emptyset$, so $V(C_1) \subseteq V(G) \setminus V(H)$.
Since $y \in C_0$, its neighbors $B_y \setminus \{q\}$ also lie in $C_0 \cup \{q\}$ in $G' - \{x, q\}$. Thus $C_1 \cap B_y = \emptyset$.
Since $B = B_x \cup B_y$, any terminal in $C_1$ must be in $B_x$. Thus $V(C_1) \setminus B = V(C_1) \setminus B_x = K$.

If $V(C_1)\nsubseteq B_x$ then $K\neq\emptyset$ immediately. Otherwise pick $b\in V(C_1)\cap B_x$.
In $G'$, the vertex $b$ has degree $3$ with neighbors $\{x\}\cup N$, where $N\subseteq V(G)\setminus V(H)$ consists of two distinct non-gadget neighbors.
After deleting $\{x,q\}$, at most one of these two neighbors can equal $q$, so there exists $w\in N\setminus\{q\}$.
Then $w\in V(C_1)$ (since $bw$ is an edge of $G'-\{x,q\}$), and $w\notin B_x$ (because $w\notin B$ and $B_x\subseteq B$).
Hence $w\in K$, so $K\neq\emptyset$.
Any edge leaving $K$ in $G$ cannot enter $V(H)$ because $K \cap B = \emptyset$. Any edge in $G$ between vertices in $V(G) \setminus V(H)$ is also in $G'$. If $v \in K$ had a neighbor $w \in G - (K \cup \{q\} \cup B_x)$, then $w$ would be a neighbor of $v$ in $G' - \{x, q\}$ outside $C_1$, which contradicts that $C_1$ is a component. Thus all neighbors of $K$ in $G$ lie in $\{q\} \cup B_x$. This proves (ii).
\end{proof}

\end{document}
